\section{Visão sistêmica e sociotécnica}
\begin{table}[!ht]\centering\small\renewcommand{\arraystretch}{1.25}
\begin{tabularx}{\linewidth}{|>{\bfseries\raggedright\arraybackslash}p{4cm}|>{\raggedright\arraybackslash}X|}
\hline
Sistema / foco & Processo de comunicação entre a escola e sua comunidade \\
\hline
Ambiente externo & Famílias, instituições de ensino, calendário escolar e eventos externos \\
\hline
Fronteira & Direção, coordenação, secretaria, professores, alunos e responsáveis \\
\hline
Entradas & Avisos, notas, frequência, calendário, mensagens e solicitações \\
\hline
Transformações & Produção, aprovação, categorização, distribuição e acompanhamento \\
\hline
Saídas & Comunicados entregues, notificações, relatórios e consultas acadêmicas \\
\hline
Dados relevantes & Usuários, acessos, avisos, notas, frequência, calendário e mensagens \\
\hline
Tecnologias atuais & WhatsApp, murais físicos e comunicados impressos \\
\hline
\end{tabularx}
\end{table}

Direção, coordenação, professores, alunos e responsáveis participam de processos que dependem do compartilhamento adequado das informações. Mudanças nos processos, dados ou tecnologias afetam a qualidade da comunicação e o engajamento da comunidade escolar.
